\chapter{Admission динамических parent-данных charged mediator}
\label{ch:v8-charged-mediator-dynamic-data-admission}

\section{Каноническое вложение старой ячейки}

Старая и расширенная ячейки имеют разложения
\begin{equation}
 \mathcal K_{43}=\mathbb C|0\rangle\oplus\mathcal N_{42},
 \qquad
 \mathcal K_{45}=\mathcal K_{43}\oplus
 \mathbb C|+\rangle\oplus\mathbb C|-\rangle.
 \label{eq:v8-mediator-admission-cell-splitting}
\end{equation}
Отсюда без выбора базиса в старом блоке возникает каноническая изометрия
\begin{equation}
 \iota:\mathcal K_{43}\hookrightarrow\mathcal K_{45},
 \qquad \iota^*\iota=I_{43},
 \qquad \rank(\iota\iota^*)=43.
 \label{eq:v8-mediator-admission-isometry}
\end{equation}
Она сохраняет опорный вакуум:
\begin{equation}
 \iota|0\rangle_{43}=|0\rangle_{45}.
 \label{eq:v8-mediator-admission-vacuum}
\end{equation}
Для локальных parents
\begin{equation}
 h_{43}=\sum_{a=1}^{42}|a\rangle\langle a|,
 \qquad
 h_{45}=h_{43}+7(P_++P_-)
 \label{eq:v8-mediator-admission-local-parents}
\end{equation}
выполнено точное compression-тождество
\begin{equation}
 \iota^*h_{45}\iota=h_{43}.
 \label{eq:v8-mediator-admission-parent-compression}
\end{equation}
Следовательно, новый product-vacuum parent является консервативным
расширением старого, а не его заменой.

\section{Real- и gauge-блок лежит в ортогональном дополнении}

Пусть
\begin{equation}
 Q_{\rm new}=I_{45}-\iota\iota^*,
 \qquad \rank Q_{\rm new}=2.
 \label{eq:v8-mediator-admission-new-projector}
\end{equation}
На его образе действуют
\begin{equation}
 Y_{45}=0_{43}\oplus\operatorname{diag}(1,-1),
 \qquad J_{45}|+\rangle=|-\rangle.
 \label{eq:v8-mediator-admission-charge-real}
\end{equation}
Поэтому старый блок gauge-тривиален относительно нового заряда и
\begin{equation}
 \iota^*Y_{45}\iota=0_{43},
 \qquad J_{45}\iota=\iota J_{43}.
 \label{eq:v8-mediator-admission-symmetry-compression}
\end{equation}
Ни одна из 42 прежних jump-меток не смешивается с новой charged Real-парой.

\section{Консервативное вложение цепи}

Покомпонентная изометрия определяет вложение опорных бесконечных тензорных
произведений
\begin{equation}
 \mathcal I=\bigotimes_{m\in\mathbb Z}\iota_m:
 (\mathcal K_{43},|0\rangle)^{\otimes\mathbb Z}longrightarrow
 (\mathcal K_{45},|0\rangle)^{\otimes\mathbb Z}.
 \label{eq:v8-mediator-admission-chain-isometry}
\end{equation}
Так как во всех узлах используется одна и та же карта, сдвиги сплетаются:
\begin{equation}
 S_{45}\mathcal I=\mathcal I S_{43}.
 \label{eq:v8-mediator-admission-shift-intertwining}
\end{equation}
При нулевом новом coupling локальный collision также сохраняет старое
подпространство, поэтому
\begin{equation}
 V_{45}(g=0)\mathcal I=\mathcal I V_{43}.
 \label{eq:v8-mediator-admission-floquet-intertwining}
\end{equation}
Тем самым все старые reduced correlators и 42-channel generator
восстанавливаются ограничением:
\begin{equation}
 \left.\mathcal L_{44}\right|_{g=0}=\mathcal L_{42}.
 \label{eq:v8-mediator-admission-generator-restriction}
\end{equation}
Различие GNVW-индексов $43$ и $45$ не противоречит этому тождеству:
$S_{43}$ является ограничением $S_{45}$ на собственную invariant
подалгебру, а не гомотопией двух полных автоморфизмов.

\section{Что действительно уже доступно}

На уровне формы ранее построены все пять требуемых типов:
\begin{equation}
 N_{\rm available\ shapes}=5/5:
 \{\text{conditional carrier},H_E,C_{1,2},S_d,
 \Gamma=\chi^2E_C/\hbar\}.
 \label{eq:v8-mediator-admission-shape-ledger}
\end{equation}
Само консервативное вложение проходит восемь независимых проверок:
\begin{equation}
 N_{\rm structural\ admission}=8/8.
 \label{eq:v8-mediator-admission-structural-ledger}
\end{equation}
Однако ни одна из форм не выбирает её физический коэффициент или статус
данных raw-parent:
\begin{equation}
 N_{\rm physical\ origin}=0/5:
 \{H_{23/24},7\gamma,g,S_{45},\tau_C\}.
 \label{eq:v8-mediator-admission-origin-ledger}
\end{equation}
Минимальный остаточный пакет новых данных тем самым локализован как
\begin{equation}
 \mathcal D_{\min}=\{\text{endpoint extension},\gamma,\chi,
 \text{transport primitive},E_C\}.
 \label{eq:v8-mediator-admission-minimal-data}
\end{equation}

\begin{theorem}[Консервативный structural admission]
Расширение $\mathcal K_{43}\subset\mathcal K_{45}$ канонически сохраняет
вакуум, локальный parent, старые 42 jump-метки, цепной сдвиг и редуцированную
динамику при нулевом новом coupling. Поэтому charged mediator совместим со
всей прежней динамической архитектурой без её перенастройки. Этот полный
structural admission не выводит пять физических входов.
\end{theorem}

\begin{nogocriterion}[Запрет превращения restriction в selector]
Тождество $S_{45}\mathcal I=\mathcal I S_{43}$ доказывает совместимость, но
не происхождение нового сдвига. Compression $\iota^*h_{45}\iota=h_{43}$ не
выбирает коэффициент charged-блока. Аналогично normalized quadratures не
фиксируют $g$, а clock-rate формула не фиксирует $E_C$ или $\chi$.
\end{nogocriterion}

\begin{protocol}\begin{enumerate}
\item каноническая изометрия $43\to45$: \textbf{есть};
\item вакуум сохраняется: \textbf{да};
\item старый parent есть compression нового: \textbf{да};
\item размерность ортогонального дополнения: \textbf{$2$};
\item дополнение Real-замкнуто: \textbf{да};
\item 42 старые jump-метки сохраняются: \textbf{да};
\item chain shifts сплетаются: \textbf{да};
\item старая reduced dynamics восстанавливается при $g=0$: \textbf{да};
\item structural admission: \textbf{$8/8$};
\item available shapes: \textbf{$5/5$};
\item physical parent-origin: \textbf{$0/5$};
\item следующий гейт: \textbf{минимальный пакет новых динамических данных}.
\end{enumerate}\end{protocol}

Машинный сертификат сохранён в
\path{s2t_v8_baryon_c0_charged_mediator_dynamic_parent_data_admission_gate_results.json}.

\begin{thebibliography}{9}
\bibitem{v8-mediator-admission-attal} S.~Attal and Y.~Pautrat,
\emph{From Repeated to Continuous Quantum Interactions},
Ann. Henri Poincar\'e 7 (2006) 59; arXiv:math-ph/0311002.
\bibitem{v8-mediator-admission-gnvw} D.~Gross, V.~Nesme, H.~Vogts and
R.~F.~Werner, \emph{Index Theory of One Dimensional Quantum Walks and
Cellular Automata}, Commun. Math. Phys. 310 (2012) 419;
arXiv:0910.3675.
\end{thebibliography}